\notechapter{N-20240330}{Combinatorial Geometry, Euler Formula, Platonic Solids, Degree, and Quotients}

\section{Good triangles}

The first problem concerns $n$ points in the plane with no three collinear.  A triangle formed by three of these points is called good if its area is maximal among all triangles formed by the point set.  The goal is to show that there cannot be more than $n$ good triangles.

The note begins with small cases.  For $n=4$, a convex quadrilateral gives four good-looking candidate triangles by deleting one vertex.  For $n=5$, a pentagon suggests five candidates.  These pictures motivate the expected upper bound.

\section{Why vertices must lie on the convex hull}

If a triangle uses a point strictly inside the convex hull, it cannot be maximal.  Moving that interior vertex outward toward a hull vertex while keeping the opposite side fixed increases area.  Therefore every good triangle has all three vertices on the boundary of the convex hull.

This is the first reduction: the problem becomes cyclic.  We may list the hull vertices in clockwise order and encode a triangle by the three arc lengths between its consecutive vertices.

\section{Cyclic arc encoding}

Let a good triangle have vertices $A,B,C$ on the convex hull, ordered cyclically.  Let
\[
  a,b,c
\]
be the numbers of boundary steps along the three arcs from $A$ to $B$, from $B$ to $C$, and from $C$ back to $A$.  Then
\[
  a+b+c=m,
\]
where $m$ is the number of hull vertices.

The handwritten argument tries to order all good triangles by their arc data.  If one arc is too large compared with another, then shifting a vertex along the circle or convex boundary can produce a triangle of larger area, contradicting maximality.  The intended conclusion is that good triangles are controlled by a cyclic balancing condition, so there can be at most one good triangle associated with each starting vertex.  Hence the number of good triangles is at most $m\le n$.

\section{The geometric idea behind the proof}

The core geometric principle is this: with one side fixed, area is proportional to the distance from the third vertex to the line containing the side.  Therefore, if a vertex is not as far as possible from the opposite side, the triangle cannot be maximal.

The diagrams in the note use this repeatedly.  If a candidate triangle sits inside a larger triangle or if one can replace a vertex by a farther hull vertex on the same side of a line, the area increases.  These local replacements are the geometric engine of the proof.

\section{Euler formula for convex polyhedra}

The next part reviews Euler's formula for a convex polyhedron:
\[
  V-E+F=2.
\]
The source gives an induction proof by contracting an edge.  Suppose a polyhedron has $v'=k+1$ vertices.  Contract one edge, merging its two endpoints into one vertex.  Then
\[
  v'-1=k,\qquad e'-1=e,\qquad f'=f.
\]
If the contracted polyhedron satisfies $v-e+f=2$, then the original one does too:
\[
  (v'-1)-(e'-1)+f'=2
  \quad\Longrightarrow\quad
  v'-e'+f'=2.
\]

The proof is schematic rather than fully formal, but it captures the invariant idea: a safe local contraction changes $V$ and $E$ together and keeps $V-E+F$ unchanged.

\section{Regular convex polyhedra}

The note then considers regular convex polyhedra.  Suppose every face is a regular $a$-gon and exactly $b$ faces meet at each vertex.  Counting incidences gives
\[
  aF=2E,\qquad bV=2E.
\]
Substitute into Euler's formula:
\[
  \frac{2E}{b}-E+\frac{2E}{a}=2.
\]
Dividing by $2E$ gives
\[
  \frac1a+\frac1b>\frac12.
\]
Since $a,b\ge3$, the only possibilities are
\[
  (a,b)=(3,3),(3,4),(4,3),(3,5),(5,3).
\]
These correspond to the tetrahedron, octahedron, cube, icosahedron, and dodecahedron.

The handwritten page lists sample triples $(F,E,V)$ such as
\[
  (4,6,4),\quad (8,12,6),\quad (6,12,8),\quad (20,30,12),\quad (12,30,20).
\]

\section{Euler characteristic}

The note marks the extension to Euler characteristic:
\[
  \chi=V-E+F.
\]
For a convex polyhedron, $\chi=2$.  More generally, $\chi$ is a topological invariant: under suitable deformations or subdivisions it remains unchanged.  This explains why Euler's formula is not merely a counting trick but a topological statement.

\section{Degree of a map}

The last page quotes the degree of a map between spheres.  For any $n>0$ and a map
\[
  f:S^n\to S^n,
\]
the induced map on top homology is multiplication by some integer $d$:
\[
  f_*:H_n(S^n)\to H_n(S^n),\qquad f_*(1)=d.
\]
This integer is the degree of $f$.

The key property is multiplicativity:
\[
  \deg(g\circ f)=\deg(g)\deg(f).
\]
The note's margin comment asks why this should be true.  The answer is that induced maps on homology compose functorially:
\[
  (g\circ f)_*=g_*\circ f_*.
\]
If $f_*$ multiplies by $d_f$ and $g_*$ multiplies by $d_g$, their composition multiplies by $d_gd_f$.

\section{Homology group and quotient group reminders}

The page also records the quotient definition
\[
  H_n(X)=Z_n(X)/B_n(X),
\]
meaning that homology is cycles modulo boundaries.  It then recalls quotient groups.  If $N\trianglelefteq G$, the quotient group $G/N$ consists of cosets, and multiplication is defined by
\[
  (g_1N)(g_2N)=(g_1g_2)N.
\]
Normality is exactly the condition needed for this product to be well-defined.

The handwritten note draws a circle and a cylinder-like picture, suggesting the intuition that quotienting identifies points and changes the topology.  Algebraic quotient and topological quotient are different constructions, but both create new objects by identifying equivalent elements.

\section{A bridge outward}

This note moves from extremal planar geometry to polyhedral topology and then to homological algebra.  The bridge is invariant thinking.  A maximal-area triangle is controlled by convex hull geometry.  A polyhedron is controlled by $V-E+F$.  A map of spheres is controlled by degree.  A quotient group is controlled by which elements are declared equivalent.  In each case, the problem becomes manageable after finding the invariant that survives the allowed transformations.

\section{Euler characteristic as invariant}

For a convex polyhedron,
\[
  V-E+F=2.
\]
This number is the Euler characteristic of the sphere.  More generally, triangulated surfaces have Euler characteristic independent of the triangulation.

This is the topological invariant behind the combinatorial formula.

\section{Degree counting on graphs}

The identity
\[
  \sum_v \deg(v)=2E
\]
counts incidences between vertices and edges.  In polyhedron problems, this turns local face/vertex constraints into global inequalities.

The classification of Platonic solids follows by combining Euler's formula with degree and face-size bounds.

\section{Quotients and surface construction}

Identifying edges of polygons constructs surfaces.  The torus, sphere, projective plane, and Klein bottle can all be represented by edge-gluing diagrams.  Euler characteristic tracks how cells survive the quotient.

\section{Euler bookkeeping from graphs to surfaces}

Combinatorial geometry becomes topology when one asks which quantities survive deformation or subdivision.  Euler characteristic, incidence counting, convex hull structure, and quotient surfaces are the advanced frame behind the finite diagrams.

\section{Euler characteristic as topological bookkeeping}

Euler's formula for a connected planar graph is
\[
  V-E+F=2.
\]
The quantity
\[
  \chi=V-E+F
\]
is the Euler characteristic.  For a triangulated surface, the same alternating count generalizes to
\[
  \chi=\sum_{k\ge0}(-1)^k f_k,
\]
where $f_k$ is the number of $k$-simplices.

The important point is invariance: if the triangulation is refined without changing the underlying surface, $\chi$ does not change.  Thus Euler's formula is not only a planar counting trick; it is a topological invariant.

\section{Euler characteristic and homology}

For a finite simplicial complex, homology groups $H_k$ measure $k$-dimensional holes.  Their dimensions
\[
  b_k=\dim H_k
\]
are Betti numbers.  The Euler characteristic also equals
\[
  \chi=\sum_{k\ge0}(-1)^k b_k.
\]
For a sphere, $b_0=1$, $b_2=1$, and all other Betti numbers vanish, so $\chi=2$.  For a torus, $b_0=1$, $b_1=2$, $b_2=1$, so $\chi=0$.

This explains why the same number appears in graph embeddings, triangulations, and surface classification.

\section{Degree theory}

The degree of a map is an integer measuring how many times the domain wraps around the target, counted with orientation.  For a map $f:S^1\to S^1$, the degree is the winding number.  If
\[
  f(e^{it})=e^{int},
\]
then $\deg f=n$.

In higher dimensions, degree theory is used to prove existence theorems.  A nonzero degree often forces a solution because the map cannot be deformed away from a target point without crossing it.

Thus the elementary ``degree'' and quotient-space discussions in the note connect to a robust topological invariant.

\section{Quotient spaces and identifications}

A quotient space is obtained by declaring some points equivalent and treating each equivalence class as one point.  For example, gluing the two ends of an interval gives a circle:
\[
  [0,1]/(0\sim1)\cong S^1.
\]
Gluing edges of a square can produce a cylinder, Möbius strip, torus, or projective plane depending on the identifications.

The elementary quotient examples are therefore not mere pictures.  They are the language by which topology builds new spaces from old ones.

\section{Euler characteristic as alternating rank}

For a chain complex
\[
  0\to C_n\xrightarrow{\partial_n}C_{n-1}\to\cdots\to C_0\to0,
\]
with finite-dimensional chain groups, the Euler characteristic satisfies
\[
  \sum_k(-1)^k\dim C_k=\sum_k(-1)^k\dim H_k.
\]
This identity is the algebraic reason a count of vertices, edges, and faces equals a count of homological holes.

For a polyhedral sphere, the chain groups count cells and the homology counts one connected component and one enclosed void.  Hence $\chi=2$.

\section{Degree and fixed-point consequences}

Degree theory is not just a way to label maps.  It proves existence.  If a continuous map $f:S^n\to S^n$ has nonzero degree, it is surjective.  Otherwise a missing point would allow the map to contract through $S^n\setminus\{p\}$, whose top homology vanishes, contradicting nonzero degree.

This is the topological form of a familiar elementary idea: a nonzero invariant prevents deformation to a trivial configuration.  Degree is one such invariant.
